\section{第四日第二讲}

在这节轻松的讲次里，我们转到整体情形，并给出一些动机与历史。我们先补上上一讲没来得及讲完的一件事：沿着我们已经走过的推理，甚至可以得到一般的无分歧 Plancherel 公式。

设 $F$ 是局部域。「数据包」$(C_c^\infty(X_F)^{G_O},\delta_e,\langle-,-\rangle)$ 于是可以完全用对偶的 $\check M$ 来描述，也就是说，Hecke 作用可以对角化。我们在一处轻微的简化假设下描述答案，即假设 $\check M$ 总形如 $\check G\times_{\check G_X}V_X$（$V_X$ 是某个 $\check G_X$-表示），且其上的 $\mathbb G_m$-作用完全来自 $V_X$ 上的作用。此时有同构
\[
  C_c^\infty(X_F)^{G_O}\overset{\sim}{\longrightarrow}\mathbb C[\check G_X]^{\check G_X},
\]
它把 $\delta_e$ 送到常函数 $1$，并把内积 $\langle-,-\rangle$ 送到积分
\[
  \langle-,-\rangle\mapsto\int_{\check G_{X,\mathrm{compact}}}\langle-,-\rangle\cdot(\text{$V_X$ 的 q-特征标})\cdot d_{Haar}.
\]
回忆 Haar 测度与 $V_X$ 的 q-特征标之积称为 Plancherel 测度。

局部情形到此为止；现在转到整体情形。我们的两个理论 $A_G$ 与 $B_{\check G}$ 将在算术流形上取值，其中 $Y$ 要么是 $\mathbb F_q$ 上的曲线，要么是代数闭包 $\overline{\mathbb F}_q$ 上的曲线。前一种情形 $A_G(Y)$ 与 $B_{\check G}(Y)$ 都是向量空间；后一种情形两者都是范畴。$G\circlearrowleft X$ 与 $\check G\circlearrowleft\check X$ 的数据（真正应当想的是 $M$ 与 $\check M$ 的数据）定义了对象 $\Theta_{G,X}\in A_G(Y)$ 与 $\Theta_{\check G,\check X}\in B_{\check G}(Y)$。本讲只专注于定义 $\Theta_{G,X}\in A_G(Y)$。（一切在数域情形同样成立，但为方便起见我们只讲函数域情形。）

我们有 $\mathbb F_q$ 上的曲线 $\Sigma$，向量空间 $A_G(\Sigma)$ 是 $G$-丛的集合 $[G]$ 上的函数空间。这些 $G$-丛定义在 $\mathbb F_q$ 上。$G$-丛集合 $[G]$ 既可以用阿代尔方式理解，也可以用几何方式理解。设 $K$ 是 $\Sigma$ 的函数域，则 $[G]$ 可以写成 $G(K)\backslash G(\mathbb A)/G(\hat O)$；$[G]$ 的点记作 $[g]$，而对每个闭点 $s\in\Sigma$，记 $g_s\in G(\hat O_s)$ 为相应的元素。

周期函数 $\Theta_{G,X}\in A_G(\Sigma)$ 可以用三种不同方式定义。

\begin{itemize}[leftmargin=1.6em]
  \item 几何方式：$\Theta_{G,X}([g])$ 是与 $[g]$ 相伴的 $X$-丛的截面个数。回忆这个 $X$-丛就是 $(X\times[g])/G$。
  \item $\Theta_{G,X}([g])$ 是满足下述条件的截面 $x\in X(K)$ 的个数：对一切闭点 $s\in\Sigma$，截面 $xg_s$ 属于 $X(\hat O_s)$。要看出这与第一条的数量相同，注意把这个 $X$-丛限制到曲线的泛点就得到 $X$（因为丛局部平凡）；于是与 $[g]$ 相伴的 $X$-丛的截面在一般纤维上就是一个点 $x\in X(K)$。要求该截面整体有定义，等价于要求 $x$ 在 $\Sigma$ 的任何闭点处都没有极点，即 $x\in X(\hat O_s)$。
  \item 算术方式：$\Theta_{G,X}([g])$ 是求和 $\sum_{x\in X(K)}\Phi(x\cdot g)$，其中 $\Phi$ 是 $\prod X(\hat O_s)$ 在 $X(\mathbb A)$ 中的示性函数。
\end{itemize}

下面算几个例子。取 $\Sigma=\mathbb P^1$、$G=\mathbb G_m$。此时 $\mathbb P^1$ 上的 $\mathbb G_m$-丛只由一个整数确定（即 $n$ 对应 $\mathcal O(n)$），故 $[G]=\mathbb Z$。再取 $X=\mathbb A^1$。取相伴的 $X$-丛无非是取该线丛的全空间，于是 $\Theta_{G,X}=:\Theta_X$ 把 $n$ 送到 $\mathcal O(n)$ 的整体截面个数。换句话说：
\[
  \Theta_X(n)=\begin{cases}q^{n+1},& n\ge0,\\ 1,& n\le-1.\end{cases}
\]
一个重要的注记是：在我们那篇论文里，需要对周期函数作正确的归一化。所以我们其实用的是周期函数的一个轻微变体，在这个例子里就是把上面的函数 $\Theta_X$ 乘上因子 $q^{-n/2}$。于是这个重标定把
\[
  n\mapsto\begin{cases}q^{n/2+1},& n\ge0,\\ q^{-n/2},& n\le-1.\end{cases}
\]
注意现在它关于 $n=-1$ 对称。（周期函数的这个归一化是由 $L^2$-理论驱动的。）

再做一个例子。设取定子群 $H\subseteq G$。此时可以取 $X=G/H$，并考虑 $\Theta_X:=\Theta_{G,X}$ 与某个函数 $f:[G]\to\mathbb C$ 的「配对」。（当 $G=\mathrm{PGL}_2$、$H=\mathbb G_m$ 时，这个配对在 1930 年代由 Hecke 算出。）所谓「配对」指的是：

\begin{itemize}[leftmargin=1.6em]
  \item 用阿代尔均匀化表述，是积分 $\int_{G(K)\backslash G(\mathbb A)}\Theta_X(g)f(g)\,dg$。这里取 Haar 测度使 $G(\hat O)$ 的体积为 $1$。
  \item 用几何图景表述，是求和 $\sum_{[g]\in[G]}\frac{1}{\#\mathrm{Aut}([g])}\Theta_X(g)f(g)$。例如在 $\mathbb G_m$ 情形，点 $[\mathcal O(n)]$ 的测度是 $1/\#\mathbb G_m(\mathbb F_q)=1/(q-1)$。
\end{itemize}

两者给出同一个数。一个好练习是验证：当 $X=G/H$ 时，这两个配对都可以用 $H$ 改写。即
\[
  \sum_{[g]\in[G]}\frac{1}{\#\mathrm{Aut}([g])}\Theta_X(g)f(g)=\sum_{[h]\in[H]}\frac{1}{\#\mathrm{Aut}([h])}f(h),
\]
\[
  \int_{G(K)\backslash G(\mathbb A)}\Theta_X(g)f(g)\,dg=\int_{[H]}f(h)\,dh,
\]
而且这些数全都相等，给出 $\langle\Theta_X,f\rangle$。这是个有帮助的练习。

下面我们来看这类配对 $\langle\Theta_X,f\rangle$ 的研究史，以及数论学家最初为何对它们感兴趣。为此一直回溯到 Riemann zeta 函数 $\zeta(s)=\sum_{n=1}^\infty n^{-s}$。它有三条性质，合起来极其刚性：

\begin{itemize}[leftmargin=1.6em]
  \item 它有 Euler 积 $\prod_p(1-p^{-s})^{-1}$。
  \item 它有解析（指亚纯）延拓。
  \item 它有函数方程，联系 $s$ 与 $1-s$ 处的值。
\end{itemize}

一个函数要同时具备这三条性质并不容易；我们就把这类函数称作「L-函数」。人们开始发现这样的 L-函数，而它们全都与算术有关：与 Dirichlet 特征标、数域、Galois 表示相联系，后来又与椭圆曲线相联系。但要证明这些函数具备上述全部性质却很难！重要的是，所有这类（预期的）L-函数都与算术本性的对象有关。

Hecke 就在这时登场了，大约是 1936 年。他发现了一种完全不同的办法，从全纯模形式（它们是如今所谓 Hecke 算子的本征函数）产生这类函数。即给定 $f(q)=\sum a_nq^n$，可以写下 L-函数 $L(f,s)=\sum a_nn^{-s}$。有 Euler 积这条性质可追溯到 Ramanujan，而它有解析延拓与函数方程的证明归功于 Hecke。他知道 Riemann 证明 zeta 函数函数方程时用到了 $\theta$-函数，这很可能是给他的重要动机。Hecke 的关键观察是：若把全纯模形式 $f$ 看作自守形式 $f:[\mathrm{PGL}_2]\to\mathbb C$，则
\[
  \int_{[\mathbb G_m]}f(h)|h|^s\,dh=L(f,s).
\]
不必在意此式中的特征标 $|h|^s$；撇开这一点，它就是 $f$ 与周期函数 $\Theta_{\mathrm{PGL}_2,\mathrm{PGL}_2/\mathbb G_m}$ 的配对。这很了不起，因为它是这类 L-函数的一个解析的（而非算术的）来源。

Hecke 之后，人们开始发现其他几种用分析方法产生 L-函数的途径。下一个例子属于 1940 年代的 Rankin 与 Selberg。他们从一对 Hecke 本征形式 $f,g$ 产生 L-函数。若写 $f(q)=\sum a_nq^n$、$g(q)=\sum b_nq^n$，则 Rankin–Selberg L-函数是 $\sum_{n\ge1}\frac{a_n\bar b_n}{n^s}$。它有 Euler 积并不太难，但它的解析延拓或函数方程则完全不显然。他们证明这些的方式，用我们的语言来说，就是把 $f\boxtimes g$（看作 $G=\mathrm{GL}_2\times\mathrm{GL}_2$ 的自守形式）与 $\Theta_{G,X}$ 配对，其中 $X=\mathrm{GL}_2\times\mathbb A^2$。这里 $\mathrm{GL}_2$ 的两个因子分别从左、右作用在 $\mathrm{GL}_2$ 上，而只有第二个 $\mathrm{GL}_2$ 作用在 $\mathbb A^2$ 上。Rankin 的动机并不是研究素数，他感兴趣的是估计 $|a_n|$。这个证明非常重要，因为它为 Deligne 证明 Weil 猜想时用到的某些关键不等式提供了启发。

此后有一段时间没有更多例子。但 Godement 与 Jacquet 在 1970 年前后，从 $\mathrm{GL}_n$ 上的 Hecke 本征形式 $f$ 出发生成了 L-函数。他们同样是靠证明该 L-函数由 $f$ 与 $\Theta_{G,X}$ 配对给出（其中 $G=\mathrm{GL}_n\times\mathrm{GL}_n$、$X=\mathrm{Mat}_n$）来证明它有解析延拓与函数方程。在所有这些论文里，你能看到的是积分以及对它们的操作；把这件事概念化为「研究某个空间」要晚得多。此后自守形式理论兴起，上述情形被大量推广。

1960 年代前后，Langlands 猜想所有 L-函数（包括所有算术 L-函数）都来自自守形式。这很重要，因为来自自守形式的 L-函数有很多例子和工具可用，而算术 L-函数几乎什么都没有。这成为发展自守形式 L-函数理论的极强动机。

我们要强调，Hecke、Rankin–Selberg、Godement–Jacquet 等人的所有这些例子都有下述形式。考虑配对 $\langle f,\Theta_{G,X}\rangle$，结果要么是 $0$，要么给出一个与 $f$ 相关的 L-函数。而配对 $\langle f,\Theta_{G,X}\rangle$ 非零的条件，恰好由对偶 Hamilton $\check G$-空间 $\check M$ 打包得很好。更详细地说，这个配对是否消没与对偶群 $\check G_X$ 有关；而配对非零时得到的 L-函数则与 $V_X$ 有关，也就是 Plancherel 测度里出现的那个表示。此外，$\Theta_{G,X}$ 只依赖余切丛 $M=T^*X$。所以大量例子由一对 $(M,\check M)$ 打包。

然而，并非所有已知例子都落在超球对偶的框架内。这是这类例子中最漂亮的一类，但最好能对框架作某种推广。
